% !TEX program = xelatex
% AmericasNLP 2026 poster — A0 portrait. PLAN-B LaTeX track.
% Primary deliverable is Lance's InDesign version; this is verified-numbers fallback + structural sandbox.
\documentclass[final]{beamer}
\usepackage[size=a0,orientation=portrait,scale=1.30]{beamerposter}
\usepackage{fontspec}
\usepackage{booktabs}
\usepackage{pgfplots}
\usepackage{qrcode}
\pgfplotsset{compat=1.17}
\usefonttheme{professionalfonts} % stop beamer substituting fonts

% --- Fonts: Charis SIL body (SLO diacritics), TeX Gyre Heros headings ---
\setmainfont{CharisSIL}[Ligatures=TeX,Extension={.ttf},UprightFont=*-Regular,ItalicFont=*-Italic,BoldFont=*-Bold,BoldItalicFont=*-BoldItalic]
\newfontfamily\headfont{texgyreheros}[Ligatures=TeX,Extension={.otf},UprightFont=*-regular,BoldFont=*-bold,ItalicFont=*-italic]
\renewcommand{\familydefault}{\rmdefault} % body in Charis SIL, not beamer's default sans
\newcommand{\lak}[1]{\textit{#1}}

% --- Palette ---
\definecolor{ink}{HTML}{1A2333}      % deep slate
\definecolor{accent}{HTML}{C0492B}   % terracotta — key emphasis
\definecolor{steel}{HTML}{2E6E8E}    % steel blue — secondary series
\definecolor{paperbg}{HTML}{F6F3EC}  % warm off-white
\definecolor{midgray}{HTML}{5B6472}
\definecolor{warmlight}{HTML}{E8C9B6} % light terracotta — for title-band accent on dark

\setbeamercolor{background canvas}{bg=paperbg}
\setbeamercolor{block title}{fg=paperbg,bg=ink}
\setbeamercolor{block body}{fg=ink,bg=white}
\setbeamerfont{block title}{family=\headfont,series=\bfseries,size=\large}
\setbeamertemplate{navigation symbols}{}
\setbeamertemplate{blocks}[rounded][shadow=false]
\setbeamertemplate{itemize items}[circle]
\setbeamercolor{itemize item}{fg=accent}
\setbeamercolor{itemize subitem}{fg=midgray}

\begin{document}
\begin{frame}[t]

% ===================== TITLE BAND =====================
\begin{beamercolorbox}[wd=\paperwidth,sep=26pt]{block title}
\centering
{\headfont\bfseries\fontsize{80}{94}\selectfont Can Frontier LLMs One-Shot Translate Lakota?\,\,{\itshape\color{warmlight}(Mostly No.)}}\\[12pt]
{\headfont\itshape\fontsize{34}{42}\selectfont Evaluating Frontier LLM Translation Capability for Lakota}\\[14pt]
{\headfont\fontsize{30}{36}\selectfont Lance Robertson\, \textperiodcentered\, University of California, San Diego}\\[4pt]
{\fontsize{24}{30}\selectfont Sixth Workshop on NLP for Indigenous Languages of the Americas (AmericasNLP 2026)}
\end{beamercolorbox}

\vspace{14pt}

% ===================== HEADLINE BAND =====================
\begin{beamercolorbox}[wd=\paperwidth,sep=24pt]{block body}
\centering
{\headfont\bfseries\fontsize{64}{76}\selectfont No frontier LLM reliably translates Lakota.}\\[14pt]
{\headfont\fontsize{40}{50}\selectfont Even measuring how badly they fail is hard.}
\end{beamercolorbox}

\vspace{20pt}

% ===================== THREE COLUMNS =====================
\begin{columns}[t]

% ---------- COLUMN 1 ----------
\begin{column}{0.31\paperwidth}

% --- Lakȟótiyapi typographic moment ---
\centering
{\rmfamily\itshape\fontsize{56}{60}\selectfont\color{accent} Lakȟótiyapi}\\[4pt]
{\fontsize{20}{24}\selectfont\color{midgray} \emph{the language under test}}
\vspace{16pt}

\begin{block}{Why test this?}
\lak{Lakȟótiyapi} (Lakota) is a Siouan language UNESCO classifies as \textbf{critically endangered} --- fewer than 2,000 first-language speakers, most over 65.
\medskip

Today's LLMs list translation support across dozens of languages, but providers don't say what they can or can't do per language. A reasonable user can't tell Lakota apart from any of the others on the list --- \emph{if it does all these, why not this one?} Direct evaluation is the only way to find out.
\end{block}

\begin{block}{Is it just memorized?}
The evaluation set comes from a published dictionary, so contamination is the obvious worry. The Lakota source appears verbatim online for ${\sim}$10\% of pairs, but the Lakota and its English reference together for only 0.5\%. Scores spread widely (6--61\% equivalence), not clustered at ceiling --- the signature of partial capability, not recall.
\end{block}

\begin{block}{Setup}
\begin{itemize}
\item \textbf{7 models.} Proprietary: Gemini 3.1 Pro, Claude Opus 4.6, Claude Sonnet 4.6, GPT-5.2. Open-weight: DeepSeek-V4-Pro, Qwen3.6-Plus, GLM-5.1.
\item \textbf{200 sentence pairs} from the New Lakota Dictionary, conversational register, \textbf{both directions}.
\item \textbf{Two conditions:} baseline vs.\ extended reasoning, where the API permits.
\item \textbf{Metrics:} chrF++ and BLEU; diacritic-normalized chrF++; an LLM \emph{semantic} judge on the English side.
\end{itemize}
\end{block}

\begin{block}{Why two judges?}
chrF++ rewards surface overlap; it cannot tell fluent-but-wrong from correct. Meaning needs a separate check. For Lakota\,$\rightarrow$\,English the model output \emph{is} English, so judging meaning takes no Lakota at all --- that's what made an LLM judge feasible.
\medskip

To keep one model from flattering its own family, every pair is judged by \textbf{two models from different families} --- Gemini 3 Flash and the independent Llama-3.3-70B --- which agree substantially: Cohen's $\kappa=0.75$ (0.89 quadratic-weighted).
\end{block}

\begin{block}{Standard Lakota Orthography}
\begin{center}{\Large \lak{ȟ}\quad\lak{ą}\quad\lak{į}\quad\lak{ų}\quad\lak{á}\quad\lak{ŋ}\quad\lak{ʼ}}\end{center}
SLO encodes phonemic distinctions with diacritics; stripping them adds \textbf{+3.3 to +6.5 chrF++} (E\,$\rightarrow$\,L) --- models get base characters right, marks wrong.
\end{block}

\end{column}

% ---------- WIDE RIGHT (chart spans 2 cols; tables + failure beneath) ----------
\begin{column}{0.64\paperwidth}

\begin{block}{Surface metrics overstate capability}
\centering
\begin{tikzpicture}
\begin{axis}[
  width=0.62\paperwidth, height=24cm,
  xmin=-0.22, xmax=1.10,
  ymin=0, ymax=70,
  xtick={0,1}, xticklabels={\textbf{chrF++}, \textbf{\% equivalent}},
  x tick label style={font=\large,yshift=-4pt},
  ytick={0,20,40,60}, y tick label style={font=\small,color=midgray},
  ylabel={\large score (Lakota\,$\rightarrow$\,English)},
  ylabel style={font=\large},
  ymajorgrids=true, grid style={dashed,gray!30,thin},
  axis lines=left,
  clip=false,
]
% Gemini — highlight: the only flat/up line (metrics agree)
\addplot[accent, very thick, mark=*, mark size=3.5pt] coordinates {(0,59.4) (1,61.3)};
\node[anchor=east, font=\normalsize\bfseries, color=accent, xshift=-6pt] at (axis cs:0,59.4) {Gemini 59.4};
\node[anchor=west, font=\normalsize\bfseries, color=accent, xshift=6pt] at (axis cs:1,61.3) {61.3};

% Sonnet (close to Opus on left — label nudged up)
\addplot[ink, thick, mark=*, mark size=2.5pt] coordinates {(0,46.1) (1,30.7)};
\node[anchor=east, font=\small\bfseries, color=ink, xshift=-6pt, yshift=6pt] at (axis cs:0,46.1) {Sonnet 46.1};
\node[anchor=west, font=\small, color=ink, xshift=6pt] at (axis cs:1,30.7) {30.7};

% Opus (close to Sonnet on left — label nudged down)
\addplot[ink, thick, mark=*, mark size=2.5pt] coordinates {(0,45.9) (1,37.0)};
\node[anchor=east, font=\small\bfseries, color=ink, xshift=-6pt, yshift=-7pt] at (axis cs:0,45.9) {Opus 45.9};
\node[anchor=west, font=\small, color=ink, xshift=6pt] at (axis cs:1,37.0) {37.0};

% DeepSeek
\addplot[ink, thick, mark=*, mark size=2.5pt] coordinates {(0,38.0) (1,19.0)};
\node[anchor=east, font=\small\bfseries, color=ink, xshift=-6pt] at (axis cs:0,38.0) {DeepSeek 38.0};
\node[anchor=west, font=\small, color=ink, xshift=6pt] at (axis cs:1,19.0) {19.0};

% GPT-5.2 (close to GLM on left — nudged up)
\addplot[ink, thick, mark=*, mark size=2.5pt] coordinates {(0,30.0) (1,13.4)};
\node[anchor=east, font=\small\bfseries, color=ink, xshift=-6pt, yshift=5pt] at (axis cs:0,30.0) {GPT-5.2 30.0};
\node[anchor=west, font=\small, color=ink, xshift=6pt, yshift=4pt] at (axis cs:1,13.4) {13.4};

% GLM (close to GPT on left — nudged down)
\addplot[ink, thick, mark=*, mark size=2.5pt] coordinates {(0,29.3) (1,9.2)};
\node[anchor=east, font=\small\bfseries, color=ink, xshift=-6pt, yshift=-6pt] at (axis cs:0,29.3) {GLM 29.3};
\node[anchor=west, font=\small, color=ink, xshift=6pt] at (axis cs:1,9.2) {9.2};

% Qwen
\addplot[ink, thick, mark=*, mark size=2.5pt] coordinates {(0,26.1) (1,7.1)};
\node[anchor=east, font=\small\bfseries, color=ink, xshift=-6pt] at (axis cs:0,26.1) {Qwen 26.1};
\node[anchor=west, font=\small, color=ink, xshift=6pt, yshift=-4pt] at (axis cs:1,7.1) {7.1};
\end{axis}
\end{tikzpicture}
\\[10pt]
{\footnotesize The picture shows the divergence; the tables below prove it. Gemini's two scores almost agree (\textcolor{accent}{terracotta}, the only flat/up line). For every other model, surface chrF++ stays in the 30s--40s while genuine meaning preservation collapses below 20\% --- chrF++ alone is misleading for low-resource translation.}
\end{block}

% --- Two tables side by side beneath the chart ---
\begin{columns}[t]
\begin{column}{0.49\textwidth}
\begin{block}{Translation quality (chrF++)}
\centering\small
\begin{tabular}{@{}lcc@{}}
\toprule
\textbf{Model} & \textbf{L\,$\rightarrow$\,E} & \textbf{E\,$\rightarrow$\,L} \\
\midrule
\multicolumn{3}{@{}l}{\itshape Proprietary}\\
Gemini 3.1 Pro & \textbf{59.4} & \textbf{42.6} \\
Claude Opus 4.6 & 45.9 & 34.3 \\
Claude Sonnet 4.6 & 46.1 & 27.0 \\
GPT-5.2 & 30.0 & 23.6 \\
\midrule
\multicolumn{3}{@{}l}{\itshape Open-weight}\\
DeepSeek-V4-Pro & 38.0 & 29.2 \\
GLM-5.1 & 29.3 & 14.6 \\
Qwen3.6-Plus & 26.1 & 18.4 \\
\bottomrule
\end{tabular}
\\[6pt]
{\footnotesize Best condition per model. Every model scores higher \emph{out of} Lakota than \emph{into} it.}
\end{block}
\end{column}
\begin{column}{0.49\textwidth}
\begin{block}{Semantic equivalence (L\,$\rightarrow$\,E)}
\centering\small
\begin{tabular}{@{}lccc@{}}
\toprule
\textbf{Model} & \textbf{Equiv.} & \textbf{Partial} & \textbf{Not} \\
\midrule
Gemini 3.1 Pro & \textbf{61.3} & 23.1 & 15.6 \\
Claude Opus 4.6 & 37.0 & 21.0 & 42.0 \\
Claude Sonnet 4.6 & 30.7 & 23.8 & 45.5 \\
DeepSeek-V4-Pro & 19.0 & 14.0 & 67.0 \\
GPT-5.2 & 13.4 & 9.8 & 76.8 \\
GLM-5.1 & 9.2 & 6.5 & 84.2 \\
Qwen3.6-Plus & 7.1 & 7.6 & 85.4 \\
\bottomrule
\end{tabular}
\\[6pt]
{\footnotesize \% of judged pairs, reasoning condition.}
\end{block}
\end{column}
\end{columns}

\vspace{6pt}

% --- Failure table | Takeaways + Origin + QR side by side ---
\begin{columns}[t]
\begin{column}{0.49\textwidth}
\begin{block}{How models fail}
\footnotesize
\begin{tabular}{@{}p{0.34\linewidth}lp{0.30\linewidth}r@{}}
\toprule
\textbf{Reference} & \textbf{Model} & \textbf{Output} & \textbf{chrF++} \\
\midrule
We live in a small trailer house with no electricity. & DeepSeek & ``small \textit{sacred} tipi'' & 18.6 \\[2pt]
Make coffee (\lak{Wakȟálya yo}). & DeepSeek & ``Sanctify it!'' & 2.9 \\[2pt]
Did my daughter call? & GLM & ``Did my younger brother come?'' & 27.6 \\[2pt]
Do you have a car? & GPT-5.2 & ``Do you have a horse?'' & 66.4 \\
\bottomrule
\end{tabular}
\\[8pt]
{\footnotesize Idioms and coined terms are translated \emph{element by element}: \lak{Wakȟálya yo} (``make coffee,'' literally ``make it boil'') becomes ``sanctify it.'' A one-word near-miss --- ``car'' rendered ``horse'' by a \emph{frontier} model --- still scores 66.4 chrF++. Surface overlap and meaning preservation diverge.}
\end{block}
\end{column}
\begin{column}{0.49\textwidth}
\begin{block}{Takeaways}
\begin{itemize}
\item No frontier LLM reliably translates Lakota, either direction (2026).
\item chrF++ overstates low-resource capability --- pair it with a meaning check.
\item Open-weight doesn't close the gap.
\end{itemize}
\smallskip
Extended reasoning changes refusal behavior more than quality (open-weight E\,$\rightarrow$\,L: \textbf{0\,$\rightarrow$\,36} refusals when enabled).
\end{block}

\begin{block}{Origin}
{\small This evaluation began as a response to a provocation from Benjamin Bratton (UCSD): \emph{``just vibe-code a Lakota translation app over the weekend.''} We built the evaluation that would test whether such an app could exist.}
\end{block}

\vspace{6pt}
\begin{beamercolorbox}[wd=\linewidth,sep=14pt]{block title}
\centering
\begin{minipage}{0.32\linewidth}\centering\qrcode[height=2.6cm,version=4]{https://github.com/robotson/lakota-translation-benchmark}\end{minipage}\hfill
\begin{minipage}{0.62\linewidth}\raggedright
{\headfont\bfseries Data \& code}\\[2pt]
{\small github.com/robotson/\\lakota-translation-benchmark}
\end{minipage}
\end{beamercolorbox}
\end{column}
\end{columns}

\end{column}
\end{columns}

\end{frame}
\end{document}
